\documentclass[11pt,a4paper]{article}

% ==============================================================================
% PAKET-PAKET PENDUKUNG LATEX
% ==============================================================================
\usepackage[utf8]{inputenc}
\usepackage[indonesian]{babel}
\usepackage{amsmath, amsfonts, amssymb}
\usepackage{geometry}
\geometry{top=2.5cm, bottom=2.5cm, left=2.5cm, right=2.5cm}
\usepackage{graphicx}
\usepackage{xcolor}
\usepackage{listings}
\usepackage{tcolorbox}
\tcbuselibrary{skins,breakable}
\usepackage{booktabs}
\usepackage{hyperref}
\usepackage{enumitem}
\usepackage{fancyhdr}
\usepackage{tikz}
\usetikzlibrary{shapes.geometric, arrows, positioning, fit, backgrounds}

% Konfigurasi Hyperref
\hypersetup{
    colorlinks=true,
    linkcolor=blue!70!black,
    citecolor=blue!70!black,
    urlcolor=blue!70!black
}

% Konfigurasi Header & Footer
\pagestyle{fancy}
\fancyhf{}
\rhead{\small \textit{Integrasi C++ NMPC USV -- Tahap 1 s.d. 4}}
\lhead{\small \textit{Panduan Pemula \& Dokumentasi Teknis}}
\rfoot{\thepage}

% Definisi Warna untuk Code Snippet
\definecolor{codegreen}{rgb}{0,0.6,0}
\definecolor{codegray}{rgb}{0.5,0.5,0.5}
\definecolor{codepurple}{rgb}{0.58,0,0.82}
\definecolor{backcolour}{rgb}{0.96,0.97,0.98}
\definecolor{darkblue}{rgb}{0.0,0.2,0.6}

\lstdefinestyle{cppstyle}{
    backgroundcolor=\color{backcolour},   
    commentstyle=\color{codegreen},
    keywordstyle=\color{blue}\bfseries,
    numberstyle=\tiny\color{codegray},
    stringstyle=\color{codepurple},
    basicstyle=\ttfamily\footnotesize,
    breakatwhitespace=false,         
    breaklines=true,                 
    captionpos=b,                    
    keepspaces=true,                 
    numbers=left,                    
    numbersep=5pt,                  
    showspaces=false,                
    showstringspaces=false,
    showtabs=false,                  
    tabsize=2,
    frame=lines
}

\lstdefinestyle{plainstyle}{
    backgroundcolor=\color{backcolour},
    basicstyle=\ttfamily\footnotesize,
    breaklines=true,
    numbers=none,
    frame=single,
    rulecolor=\color{gray!40}
}
\lstset{style=cppstyle}

% Box Kustom untuk Analogi & Catatan Penting
\newtcolorbox{analogybox}[1]{
    colback=blue!5!white,
    colframe=blue!75!black,
    fonttitle=\bfseries,
    title=#1,
    arc=4mm,
    boxrule=1pt,
    breakable
}

\newtcolorbox{notebox}[1]{
    colback=green!5!white,
    colframe=green!60!black,
    fonttitle=\bfseries,
    title=#1,
    arc=4mm,
    boxrule=1pt,
    breakable
}

\newtcolorbox{warnbox}[1]{
    colback=orange!5!white,
    colframe=orange!80!black,
    fonttitle=\bfseries,
    title=#1,
    arc=4mm,
    boxrule=1pt,
    breakable
}

% ==============================================================================
% JUDUL DOKUMEN
% ==============================================================================
\title{
    \vspace{-1.0cm}
    \Huge \textbf{Penjelasan Lengkap \& Panduan Pemula:}\\[0.3cm]
    \LARGE Integrasi Algoritma NMPC Kapal Otonom ke dalam Bahasa C++ (\texttt{CPP\_Gabungan})\\[0.2cm]
    \large \textit{Dokumentasi Teknis: Tahap 1 (Solver NMPC), Tahap 2 (Guidance \& ENU), Tahap 3 (State Observer), \& Tahap 4 (Loop Utama main.cpp)}
}
\author{\textbf{Tim Riset Kapal Autonomous (USV)}}
\date{\today}

\begin{document}

\maketitle

\begin{abstract}
Dokumen ini disusun khusus sebagai jembatan pemahaman bagi peneliti atau mahasiswa yang telah terbiasa menggunakan MATLAB namun \textbf{belum familiar dengan pemrograman bahasa C++}. Di dalamnya dijelaskan secara mendalam, konseptual, dan langkah-demi-langkah mengenai perkembangan integrasi perangkat lunak pada repositori \texttt{CPP\_Gabungan}, mencakup: \textbf{Tahap 1} (Penyatuan 124 File Codegen Solver NMPC \& Benchmark Waktu), \textbf{Tahap 2} (Modul Guidance \& Transformasi Koordinat ENU Real-Time), \textbf{Tahap 3} (State Observer, Kalibrasi Offset Yaw $+180^\circ$, dan Normalisasi State Non-Dimensional $5 \times 1$), serta \textbf{Tahap 4} (Modifikasi Loop Utama \texttt{src/main.cpp}, Integrasi Kendali Real-Time 10 Hz, Mode CLI \texttt{--rudder-mode nmpc}, dan Mekanisme Fail-Safe).
\end{abstract}

\tableofcontents
\newpage

% ==============================================================================
% BAB 1: PENGENALAN C++ UNTUK PENGGUNA MATLAB/PYTHON
% ==============================================================================
\section{Pengenalan C++: Apa Bedanya dengan MATLAB dan Python?}

Bagi Anda yang terbiasa menulis kode di MATLAB atau Python, berpindah ke C++ seringkali membingungkan karena konsepnya sangat berbeda. Mari kita pahami perbedaan mendasar ini dengan analogi sederhana.

\subsection{Bahasa Interpreter vs Bahasa Kompilasi (Compiled Language)}

\begin{itemize}[leftmargin=*]
    \item \textbf{MATLAB dan Python (Interpreted):}
    Bekerja seperti \textit{penerjemah langsung (interpreter)}. Setiap kali Anda menekan tombol \textit{Run}, ada program penerjemah yang membaca baris demi baris perintah Anda, menerjemahkannya ke bahasa mesin (0 dan 1), lalu mengeksekusinya saat itu juga. 
    \begin{itemize}
        \item \textit{Kelebihan:} Mudah ditulis, fleksibel, mudah diuji coba secara interaktif.
        \item \textit{Kekurangan:} Lambat, boros memori CPU, dan memiliki jeda waktu (\textit{latency}) yang tidak menentu.
    \end{itemize}

    \item \textbf{Bahasa C++ (Compiled):}
    Bekerja seperti \textit{koki yang memasak seluruh bahan makanan sampai matang terlebih dahulu sebelum disajikan}. Sebelum program C++ bisa dijalankan di komputer kapal (Mini PC), seluruh kode sumber harus diproses oleh sebuah alat bernama \textbf{Compiler} (misalnya \texttt{g++.exe} atau \texttt{MSVC}). Compiler mengubah seluruh kode menjadi satu file aplikasi yang langsung dimengerti prosesor komputer (\texttt{.exe} di Windows atau file binary di Linux).
    \begin{itemize}
        \item \textit{Kelebihan:} \textbf{Sangat cepat (bisa 50--100 kali lebih cepat dari MATLAB)}, sangat hemat daya, dan waktu eksekusinya sangat pasti (\textit{deterministic}).
        \item \textit{Kekurangan:} Harus dikompilasi terlebih dahulu, dan struktur file harus sangat rapi dan tertata.
    \end{itemize}
\end{itemize}

\begin{analogybox}{Analogi Sederhana: Memasak Sendiri vs Membeli Makanan Instan}
Menjalankan simulasi di MATLAB seperti memasak di dapur mewah: Anda punya banyak alat bantu otomatis, tetapi proses memasaknya butuh waktu lama dan memakan banyak tempat di meja. 

Mengubahnya ke C++ seperti mengemas resep tersebut ke dalam \textbf{ransum makanan siap saji yang super cepat dibuka di lapangan}. Di atas kapal otonom nyata, kita tidak bisa membawa seluruh software MATLAB (yang ukurannya gigabyte dan berat untuk CPU Mini PC), melainkan kita hanya membawa file kecil \texttt{read\_write\_serial.exe} hasil kompilasi C++ yang langsung bekerja secepat kilat.
\end{analogybox}

\subsection{Mengenal Pasangan File C++: Header (\texttt{.h} / \texttt{.hpp}) dan Source (\texttt{.cpp})}

Di MATLAB, satu fungsi biasanya ditulis dalam satu file \texttt{.m}. Di C++, kode hampir selalu dipisah menjadi \textbf{dua jenis file}:

\begin{enumerate}[leftmargin=*]
    \item \textbf{File Header (\texttt{.h} atau \texttt{.hpp}): ``Daftar Menu Restoran''}
    File ini hanya berisi \textit{deklarasi} (daftar nama fungsi, variabel, dan parameter apa saja yang ada), tanpa menuliskan rumus lengkapnya. Tujuannya memberi tahu file lain: \textit{``Hei, saya punya fungsi kendali NMPC dengan nama \texttt{nmpc\_kapal\_waypoint}, inputnya ada 5 array, dan outputnya sudut kemudi''}.
    
    \item \textbf{File Source (\texttt{.cpp}): ``Dapur \& Resep Memasak''}
    File ini berisi \textit{implementasi nyata} (rumus matematika, perulangan for-loop, perhitungan matriks QP, SQP solver). Di sinilah seluruh baris perhitungan algoritma dikerjakan.
\end{enumerate}

% ==============================================================================
% BAB 2: ARSITEKTUR INTEGRASI & STRUKTUR FOLDER CPP_GABUNGAN
% ==============================================================================
\section{Arsitektur Integrasi Proyek \texttt{CPP\_Gabungan}}

Struktur folder \texttt{CPP\_Gabungan} dibangun sebagai kesatuan utuh aplikasi kendali kapal autonomous:

\begin{lstlisting}[style=plainstyle, caption={Struktur Direktori Proyek CPP\_Gabungan Terkini}]
CPP_Gabungan/
|-- CMakeLists.txt              <-- "Resep" kompilasi otomatis (Build System)
|-- README.md                   <-- Dokumentasi teknis lengkap proyek
|-- start_read_write_serial.bat <-- Skrip autorun Windows Mini PC di kapal
|-- include/                    <-- Kumpulan File Header (Daftar Menu Fungsi)
|   |-- guidance.hpp            <-- [TAHAP 2] Modul Guidance, Proyeksi ENU & Horizon
|   |-- state_observer.hpp      <-- [TAHAP 3] Estimator State & Kalibrasi Yaw (+180 deg)
|   |-- serial_port.hpp         <-- Driver komunikasi USB COM Port
|   |-- telemetry_parser.hpp    <-- Parser data telemetri 8 kolom dari ESP32
|   +-- processor.hpp           <-- Operasi logika telemetri
|-- src/                        <-- Kumpulan File Source (Implementasi Rumus)
|   |-- guidance.cpp            <-- [TAHAP 2] Rumus WGS84->ENU, r_tran & Horizon
|   |-- state_observer.cpp      <-- [TAHAP 3] Normalisasi State Non-Dim [v',r',x',y',psi]
|   |-- main.cpp                <-- [TAHAP 4] Loop utama kontroler NMPC & serial bridge
|   +-- serial_port.cpp         <-- Implementasi driver serial port Windows/Linux
|-- nmpc/                       <-- Kumpulan 124 File Solver NMPC (MATLAB Codegen)
|   |-- nmpc_kapal_waypoint.h   <-- Header utama solver NMPC WyNDA
|   |-- nmpc_kapal_waypoint.cpp <-- Algoritma SQP & Model WyNDA 11 Parameter
|   |-- fmincon.cpp             <-- Non-linear optimization solver
|   +-- ... (56 file algoritma optimasi kuadratik lainnya)
+-- tests/                      <-- Modul Pengujian Mandiri (Unit Testing)
    |-- test_nmpc.cpp           <-- [TAHAP 1] Pengujian benchmark solver NMPC
    |-- test_guidance.cpp       <-- [TAHAP 2] Pengujian guidance & konversi ENU
    +-- test_state_observer.cpp <-- [TAHAP 3] Pengujian state observer & kalibrasi yaw
\end{lstlisting}

% ==============================================================================
% BAB 3: REKAPITULASI TAHAP 1 (SOLVER NMPC & BENCHMARK)
% ==============================================================================
\section{Rekapitulasi Tahap 1: Integrasi Library Codegen \& Benchmark}

Pada Tahap 1, 124 file hasil generate MATLAB Coder berhasil diintegrasikan menjadi sebuah \textbf{Static Library} bernama \texttt{libnmpc\_solver.a} menggunakan CMake dengan flag optimasi \texttt{-O3} dan \texttt{-fopenmp}.

Hasil uji coba \texttt{test\_nmpc.exe} menunjukkan bahwa solver SQP NMPC WyNDA membutuhkan waktu rata-rata hanya \textbf{0.82 milidetik per langkah}. Artinya, pada sistem kendali 10 Hz ($100\text{ ms}$ per siklus), komputer Mini PC memiliki **margin waktu luang 99.18\%**, sehingga beban CPU sangat ringan.

% ==============================================================================
% BAB 4: REKAPITULASI TAHAP 2 (GUIDANCE & TRANSFORMASI KOORDINAT ENU)
% ==============================================================================
\section{Rekapitulasi Tahap 2: Manajemen Waypoint \& Koordinat ENU}

Tahap 2 berfokus pada pembangunan **Modul Guidance** (\texttt{guidance.hpp} dan \texttt{guidance.cpp}) sebagai pemandu rute navigasi.

\subsection{Rumus Matematis Proyeksi WGS84 Geodetik ke ENU}
Diberikan titik acuan Home Point $(\text{lat}_0, \text{lon}_0)$ dan titik sembarang $(\text{lat}, \text{lon})$:
\begin{align}
    \Delta \text{Lat} &= (\text{lat} - \text{lat}_0) \times \frac{\pi}{180}, \quad \Delta \text{Lon} = (\text{lon} - \text{lon}_0) \times \frac{\pi}{180} \\
    Y_{\text{North}} &= R_{\text{earth}} \times \Delta \text{Lat}, \quad X_{\text{East}} = R_{\text{earth}} \times \cos\left(\text{lat}_0 \times \frac{\pi}{180}\right) \times \Delta \text{Lon}
\end{align}
di mana jari-jari bumi rata-rata $R_{\text{earth}} = 6,371,000\text{ meter}$.

Setiap 100 milidetik (10 Hz), posisi GPS kapal saat ini $(\text{lat}_{\text{ship}}, \text{lon}_{\text{ship}})$ dikonversi secara real-time menjadi $(X_{\text{ship}}, Y_{\text{ship}})$ pada koordinat ENU meter sebagai titik tolak pembentukan 20 langkah prediksi lintasan horizon NMPC ($x_{\text{ref\_seq}}$ dan $y_{\text{ref\_seq}}$).

% ==============================================================================
% BAB 5: REKAPITULASI TAHAP 3 (STATE OBSERVER & KALIBRASI YAW)
% ==============================================================================
\section{Rekapitulasi Tahap 3: State Observer \& Kalibrasi Yaw $+180^\circ$}

Tahap 3 bertindak sebagai penerjemah sensor telemetri ESP32-S3:
\begin{enumerate}
    \item \textbf{Kalibrasi Orientasi Yaw:} Mengoreksi pergeseran sudut fisik sensor IMU sebesar $+180^\circ$ ($\text{yaw}_{\text{cal}} = (\text{yaw}_{\text{raw}} + 180^\circ) \pmod{360^\circ}$) dan mengonversinya ke sistem ENU ($\psi = \text{wrap\_to\_pi}(\text{yaw}_{\text{cal}} \cdot \pi/180)$).
    \item \textbf{Vektor State Non-Dimensional WyNDA ($5 \times 1$):}
    \begin{equation}
        s = \begin{bmatrix} v' \\ r' \\ x' \\ y' \\ \psi \end{bmatrix} = \begin{bmatrix} v_{\text{sway}} / u_0 \\ (\text{yaw\_rate}_{\text{dps}} \cdot \frac{\pi}{180}) \cdot \frac{L}{u_0} \\ X_{\text{ENU}} / L \\ Y_{\text{ENU}} / L \\ \psi_{\text{ENU}} \end{bmatrix}
    \end{equation}
    *(dengan $L = 1.0107\text{ m}$ dan $u_0 = 0.6114\text{ m/s}$)*.
\end{enumerate}

% ==============================================================================
% BAB 6: PENJELASAN MENDALAM TAHAP 4 (MODIFIKASI LOOP UTAMA MAIN.CPP)
% ==============================================================================
\section{Penjelasan Mendalam Tahap 4: Integrasi Penuh Loop Utama \texttt{main.cpp}}

Tahap 4 adalah tahap penyatuan di mana seluruh modul yang telah dibangun (Serial Communication, State Observer, Guidance Module, dan SQP NMPC Solver) dirangkai ke dalam program eksekutabel utama \texttt{read\_write\_serial.exe}.

\subsection{Arsitektur Siklus Kendali Real-Time 10 Hz}

Kapal beroperasi dengan siklus data 10 Hz (satu kali pertukaran data setiap 100 milidetik). Gambar~\ref{fig:loop_nmpc} mengilustrasikan alur pemrosesan di dalam loop utama \texttt{main.cpp}:

\begin{figure}[htbp]
\centering
\begin{tikzpicture}[
    block/.style={
        rectangle, 
        draw=darkblue, 
        fill=blue!5!white, 
        text width=8.5cm, 
        align=left, 
        rounded corners=4pt, 
        line width=0.9pt,
        inner sep=6pt,
        font=\footnotesize
    },
    arrow/.style={
        draw=darkblue, 
        -latex', 
        line width=1.2pt
    }
]
    \node [block] (step1) {\textbf{1. Serial RX:} Menerima baris CSV telemetri 8 kolom dari ESP32};
    \node [block, below=0.35cm of step1] (step2) {\textbf{2. Tag [WP] Handler:} Update Home \& Waypoints jika ada rute baru dari darat};
    \node [block, below=0.35cm of step2] (step3) {\textbf{3. State Observer:} Kalibrasi Yaw $+180^\circ$, hitung $(X_{\text{ship}}, Y_{\text{ship}})$, buat $s[5]$};
    \node [block, below=0.35cm of step3] (step4) {\textbf{4. Guidance:} Cek $r_{\text{tran}} \le 3.0\text{ m}$, switching WP, buat horizon $N=20$};
    \node [block, below=0.35cm of step4] (step5) {\textbf{5. Solver NMPC SQP:} Hitung sudut kemudi optimal $u_{\text{opt}}$};
    \node [block, below=0.35cm of step5] (step6) {\textbf{6. Fail-Safe \& Clamping:} Batasi kemudi $[-35^\circ, +35^\circ]$, fallback jika exitflag $\le 0$};
    \node [block, below=0.35cm of step6] (step7) {\textbf{7. Serial TX:} Kirim string \texttt{timestamp,rudder\_deg} balik ke ESP32};

    \path [arrow] (step1) -- (step2);
    \path [arrow] (step2) -- (step3);
    \path [arrow] (step3) -- (step4);
    \path [arrow] (step4) -- (step5);
    \path [arrow] (step5) -- (step6);
    \path [arrow] (step6) -- (step7);
\end{tikzpicture}
\caption{Alur Siklus Kendali Real-Time 10 Hz pada \texttt{src/main.cpp}.}
\label{fig:loop_nmpc}
\end{figure}

\subsection{Pembedahan Kode Loop Utama (\texttt{src/main.cpp})}

Berikut adalah potongan inti implementasi loop NMPC pada \texttt{src/main.cpp}:

\begin{lstlisting}[language=C++]
if (rudder_mode == "nmpc") {
  // 1. Update State Observer (Kalibrasi Yaw +180 deg & Posisi ENU)
  usv::ShipState ship_state = observer.update(*row, guidance, u_prev_rad);

  if (guidance.has_home() && !guidance.get_waypoints_enu().empty()) {
    // 2. Update Guidance & Generator Horizon NMPC
    auto g_out = guidance.update_guidance(ship_state.x_enu, ship_state.y_enu);

    if (g_out.is_mission_completed) {
      rudder_cmd_deg = 0.0; // Misi selesai -> kemudi lurus aman
      u_prev_rad = 0.0;
    } else {
      // 3. Eksekusi NMPC WyNDA SQP Solver
      double u_opt_rad = 0.0;
      double exitflag = 0.0;
      nmpc_kapal_waypoint(ship_state.state_nd, u_prev_rad, surge_speed,
                          g_out.x_ref_seq, g_out.y_ref_seq, g_out.psi_ref_seq,
                          &u_opt_rad, &exitflag);

      // 4. Mekanisme Fail-Safe
      if (exitflag <= 0.0) {
        rudder_cmd_deg = u_prev_rad * 180.0 / M_PI; // Tahan kemudi sebelumnya
      } else {
        const double u_opt_deg = u_opt_rad * 180.0 / M_PI;
        rudder_cmd_deg = clamp_rudder_deg(u_opt_deg, -35.0, 35.0);
        u_prev_rad = rudder_cmd_deg * M_PI / 180.0;
      }
    }
  } else {
    rudder_cmd_deg = 0.0; // Menunggu rute waypoint dari operator
  }
}
\end{lstlisting}

\subsection{Dukungan Argumen Baris Perintah (CLI) untuk Uji Coba}

Untuk memudahkan riset dan eksperimen di kolam uji maupun danau, program mendukung berbagai opsi CLI:

\begin{lstlisting}[style=plainstyle]
# 1. Menjalankan kendali otonom NMPC penuh:
./read_write_serial.exe --port COM16 --rudder-mode nmpc --print nmpc

# 2. Uji coba dengan 4 Waypoint default tanpa dashboard darat:
./read_write_serial.exe --port COM16 --rudder-mode nmpc --default-test-wp --print all

# 3. Menyesuaikan parameter fisik kapal secara dinamis:
./read_write_serial.exe --port COM16 --yaw-offset 180.0 --r-tran 2.5 --surge-speed 0.70
\end{lstlisting}

Selain itu, skrip autorun \texttt{start\_read\_write\_serial.bat} telah diperbarui agar Mini PC di kapal dapat otomatis mengeksekusi kendali NMPC saat pertama kali dinyalakan.

% ==============================================================================
% BAB 7: GLOSARIUM ISTILAH TEKNIS (UNTUK PEMULA)
% ==============================================================================
\section{Kamus Singkat Istilah Teknis (Glosarium)}

\begin{description}[leftmargin=*, style=nextline]
    \item[\textbf{Fail-Safe}]
    Mekanisme pengaman otomatis yang mengambil alih kontrol saat terjadi kegagalan (misalnya solver tidak konvergen atau GPS terputus) agar kapal tidak hilang kendali.

    \item[\textbf{Clamping}]
    Pembatasan nilai numerik pada batas atas dan batas bawah (misalnya membatasi kemudi di $[-35^\circ, +35^\circ]$ agar motor servo tidak rusak).

    \item[\textbf{Heartbeat (\$HB)}]
    Sinyal periodik (1 detik sekali) yang dikirim Mini PC ke ESP32 sebagai tanda bahwa komputer kendali aktif dan sehat.

    \item[\textbf{Command Line Interface (CLI)}]
    Antarmuka baris perintah teks yang memungkinkan pengguna mengatur parameter program saat program pertama kali dipanggil melalui terminal.
\end{description}

% ==============================================================================
% BAB 8: LANGKAH SELANJUTNYA (MENUJU TAHAP 5)
% ==============================================================================
\section{Rencana Tahapan Selanjutnya (Tahap 5)}

Dengan tuntasnya **Tahap 1, 2, 3, dan 4**, seluruh sistem kendali NMPC di kapal telah terintegrasi secara utuh dan siap jalan. Langkah berikutnya adalah **Tahap 5: Data Logger CSV Komprehensif**:

\begin{enumerate}
    \item \textbf{Perekaman Log Otomatis (\texttt{logs/nmpc\_log\_YYYYMMDD\_HHMMSS.csv}):}
    Menyimpan seluruh variabel telemetri, koordinat ENU, sudut target, error pelacakan (*cross-track error* \& *heading error*), serta waktu komputasi SQP setiap 100 ms.
    
    \item \textbf{Persiapan Analisis Data Jurnal / Tesis:}
    Membuat skrip plotting otomatis di Python/MATLAB untuk menghasilkan grafik lintasan trajektori 2D, respon sudut kemudi, dan kestabilan yaw kapal.
\end{enumerate}

---

\vspace{0.5cm}
\begin{center}
\textbf{--- Akhir Dokumen Penjelasan Tahap 1 s.d. 4 ---}\\
\textit{Disimpan pada berkas: \texttt{USV-NMPC/penjelasan CPP\_Gabungan/penjelasan.tex}}
\end{center}

\end{document}
